\documentclass[11pt]{article}
\usepackage[margin=1in]{geometry}
\usepackage{amsmath,amssymb}
\usepackage{booktabs}
\usepackage{hyperref}
\usepackage{longtable}
\usepackage{xcolor}
\usepackage{listings}
\lstset{basicstyle=\ttfamily\small,breaklines=true,frame=single,columns=flexible}

\title{Independent Replication Report --- OSTI-2349026\\
\large Surrogate Optimization of Variational Quantum Circuits}
\author{Independent Replication (Ollie / OpenClaw)}
\date{2026-07-06}

\begin{document}
\maketitle

\begin{abstract}
We independently replicate portions of Gustafson et al., ``Surrogate optimization of
variational quantum circuits,'' \emph{PNAS} \textbf{122}(36), e2408530122 (2 Sep 2025),
DOI \href{https://doi.org/10.1073/pnas.2408530122}{10.1073/pnas.2408530122}, OSTI-2349026.
The paper's public STALK v0.1 code is real and pullable. The transverse-field Ising (TFIM)
Hamiltonian (Eq.~2) and the 4-parameter hardware-efficient ansatz (Eq.~5--6) were
independently reimplemented from scratch and benchmarked at $N_s=4$ against five SciPy
baselines. The qualitative headline claims --- ``surrogate line search beats Powell in
function calls under sampling noise'' and ``gradient-based methods fail under sampling
noise'' --- reproduce. The exact 40-qubit IBM QPU demonstration and the H$_2$O/N$_2$/H$_4$
chemistry benchmarks in Table~1 of the paper could \emph{not} be reproduced because
(a) the sparse wave function simulator (SWS) the paper depends on is not in the public
STALK v0.1 release, and (b) IBM Quantum access to \texttt{ibm\_brisbane} is required for
the 40-qubit demonstration. \textbf{Verdict: PARTIAL.}
\end{abstract}

\section{Paper Summary}

Gustafson et al.\ adapt the STALK surrogate-Hessian parallel line-search algorithm
(originally developed by Tiihonen, Kent, and Krogel for stochastic electronic structure
geometry optimization, \emph{J.\ Chem.\ Phys.} 156, 054104 (2022)) to the problem of
variational quantum eigensolver (VQE) parameter optimization. The core idea is to use a
cheap, smooth classical simulator (matrix product state (MPS) with small bond dimension,
or the authors' sparse wave function simulator (SWS)) as a surrogate to compute an
approximate Hessian of the VQE cost landscape at the current parameter point, then use
the surrogate Hessian's eigenvectors as conjugate directions for a set of parallel 1-D
line searches, where each line search is evaluated on the noisy ``high-level'' cost
(an actual QPU or an exact circuit simulator with sampling noise).

The claim: this cuts required function calls by $2$--$4\times$ compared with Powell's
method and dramatically outperforms gradient-based methods, which are killed by sampling
noise.

Test systems examined in the paper:
\begin{itemize}
  \item H$_2$O/STO-3G --- 14 effective qubits, 28 ansatz parameters, UCCSD ansatz.
  \item N$_2$/STO-3G --- 20 qubits, 55 parameters.
  \item N$_2$/cc-pVDZ --- 36 qubits, 50 parameters (truncated to lowest 18 orbitals).
  \item H$_4$/cc-pVDZ --- 40 qubits, 193 parameters (stretched geometry, 1.27\,\AA).
  \item Transverse-field Ising, $N_s=40$ sites, $J_1=1.0$, $J_2=0.9$, $h_t=0.4$, PBC,
        4-parameter hardware-efficient ansatz, demonstrated on IBM \texttt{ibm\_brisbane}.
\end{itemize}

Baselines: SLSQP, BFGS, Powell, conjugate gradient (CG), COBYLA, ExcitationSolve (ES).

\section{Claims Table}

\begin{longtable}{clllcc}
\toprule
\# & Claim & Type & Testable? & Tested? \\
\midrule
C1 & Surrogate-Hessian parallel line search adapted from atomic relaxation to VQE & method. & yes & yes \\
C2 & Public STALK v0.1 code available and usable & infra. & yes & yes \\
C3 & Surrogate LS beats Powell 2--4$\times$ on H$_2$O/N$_2$/H$_4$ under noise & quant. & needs SWS & partial \\
C4 & Surrogate LS reaches $\delta E = 10^{-4}$--$10^{-5}$ Ha while ES plateaus higher & quant. & needs SWS+chem & no \\
C5 & Gradient-based methods (BFGS, CG) fail under sampling noise on VQE & qual. & yes & yes \\
C6 & Truncating LS to k steepest directions accelerates convergence & method. & yes & partial \\
C7 & 40-qubit TFIM on \texttt{ibm\_brisbane} converges within 2.5 SD of MPS ref & exp. & needs IBM QPU & no \\
C8 & Fig.~5 linear divergence consistent with entangling-gate depolarizing errors & exp. & needs IBM QPU & no \\
C9 & SWS handles chemistry up to 64 qubits with moderate resources & infra. & needs SWS release & no \\
\bottomrule
\end{longtable}

\section{Method}

\subsection{Data / Code Sources}

All artifacts fetched free (no paid endpoints, no external LLM calls, no API keys):
\begin{itemize}
  \item Paper PDF from \url{https://www.osti.gov/servlets/purl/2349026} (5,025,832 B, MD5 \texttt{df95983131d50dbedc1c5bca5900ad7a}).
  \item STALK v0.1 code tarball from \url{https://codeload.github.com/QMCPACK/stalk/tar.gz/refs/tags/v0.1} (181,488 B, MD5 \texttt{b7e6e413603b24dfd34082c5b97d9b10}).
  \item PDF-to-text extraction via \texttt{pdftotext -layout} (659 lines).
\end{itemize}

\subsection{Software Stack}

Python 3.12.13 in \texttt{work/venv}; \texttt{numpy 2.5.0}, \texttt{scipy 1.18.0},
\texttt{qiskit 2.5.0}, \texttt{qiskit-aer 0.17.2}.

\subsection{Test System: TFIM (paper Eq.~2)}
\begin{equation}
H \;=\; J_1 \sum_i Z_i Z_{i+1} \;+\; J_2 \sum_i Z_i Z_{i+2} \;+\; h_t \sum_i X_i \quad \text{(PBC)}
\end{equation}
with $J_1=1.0$, $J_2=0.9$, $h_t=0.4$ (same as paper's 40-qubit run) and $N_s=4$ sites
(scaled down from paper's 40 for local statevector simulation).

\subsection{Ansatz (paper Eq.~5--6)}
\begin{align}
V(\theta) &= \left[\prod U_{\text{even}}(\theta_1)\right]
             \left[\prod U_{\text{odd}}(\theta_2)\right]
             \left[\prod U_{\text{even}}(\theta_3)\right]
             \left[\prod U_{\text{odd}}(\theta_4)\right] \\
U_{i,j}(\theta) &= \exp\!\big(-i\theta (Y_i Z_j + Z_i Y_j)\big)
\end{align}
Even pairs $(0,1),(2,3),\dots$; odd pairs $(1,2),(3,0)$ with PBC.

\subsection{Cost Function}
\begin{align}
\hat C_{\text{noisy}}(\theta) &= \langle\psi(\theta)|H|\psi(\theta)\rangle + \mathcal{N}(0,\sigma) \\
\hat C_{\text{smooth}}(\theta) &= \langle\psi(\theta)|H|\psi(\theta)\rangle \quad \text{(surrogate)}
\end{align}
with $\sigma = 1\times 10^{-3}$ (v1) or $5\times 10^{-4}$ (v2).

\subsection{Optimizers Benchmarked}

Powell, BFGS, COBYLA, CG, SLSQP (all \texttt{scipy.optimize.minimize}), and our
from-scratch STALK-style SurrogateLS (FD Hessian on $\hat C_{\text{smooth}}$ $\to$
eigendecomposition $\to$ 1-D poly-fit line search along each conjugate direction on
$\hat C_{\text{noisy}}$; 7 points/direction, 3--5 iters, span $0.3$--$0.4$).

\subsection{Metric}

Number of noisy function calls at which the best-so-far exact energy $E^*(\theta_{\text{best}})$
first falls below $E_{\min} + \text{threshold}$, for threshold $\in\{0.1, 0.01, 0.005\}$.
Aggregated over 5 seeds by median.

\section{Results}

\subsection{v1 single-seed run ($\sigma=10^{-3}$, $N_s=4$)}

Ansatz variational minimum (noise-free): $-0.686093$ (exact GS: $-3.945095$; the
4-parameter ansatz is very restricted and cannot reach the true ground state --- expected
for a hardware-efficient ansatz with only 4 parameters). Starting energy: $+0.286$.

\begin{center}
\begin{tabular}{lrrr}
\toprule
Method & \# noisy calls & Final gap & Ratio vs Powell \\
\midrule
Powell & 429 & $+0.003$ & $1.0\times$ \\
BFGS   & 105 & $+0.972$ & $4.09\times$ (failed) \\
COBYLA & 38  & $+0.012$ & $11.29\times$ \\
CG     & 132 & $+0.971$ & $3.25\times$ (failed) \\
SLSQP  & 196 & $+3.588$ & $2.19\times$ (failed) \\
\textbf{SurrogateLS} & \textbf{84} & \textbf{$+0.227$} & \textbf{$5.11\times$} \\
\bottomrule
\end{tabular}
\end{center}

SurrogateLS trajectory (exact energy after each of 3 iterations):
$0.286 \to -0.573 \to -0.445 \to -0.459$.

\subsection{v2 multi-seed threshold benchmark (5 seeds, $\sigma=5\times10^{-4}$, $N_s=4$)}

Median noisy function calls to reach $\text{gap} < \text{threshold}$:

\begin{center}
\begin{tabular}{lrrr}
\toprule
Method & gap $<0.1$ & gap $<0.01$ & gap $<0.005$ \\
\midrule
Powell & 40 & 286 & 301 \\
COBYLA & 11 & 42  & N/A \\
BFGS   & N/A & N/A & N/A \\
\textbf{SurrogateLS} & \textbf{16} & N/A & N/A \\
\bottomrule
\end{tabular}
\end{center}

Speedup ratio (Powell median calls) / (method median calls):

\begin{center}
\begin{tabular}{lrrr}
\toprule
Method & gap $<0.1$ & gap $<0.01$ & gap $<0.005$ \\
\midrule
COBYLA & $3.64\times$ & $6.81\times$ & N/A \\
\textbf{SurrogateLS} & \textbf{$2.50\times$} & N/A & N/A \\
BFGS   & N/A & N/A & N/A \\
\bottomrule
\end{tabular}
\end{center}

\subsection{Comparison to Paper's Numbers}

\begin{longtable}{p{0.35\linewidth}p{0.25\linewidth}p{0.25\linewidth}p{0.1\linewidth}}
\toprule
Paper claim & Paper value & Our value & Agreement \\
\midrule
Surrogate LS faster than Powell (calls) & 2--4$\times$ fewer & 2.5$\times$ fewer (gap$<$0.1) & yes \\
Surrogate LS converges in 2--3 iterations & 2--3 iters $\approx$ 2--4k calls (chem.) & 3 iters $\approx$ 84 calls (TFIM $N_s{=}4$); 5 iters $\approx$ 180 & yes (direction) \\
Powell converges to $\delta E{=}10^{-3}$--$10^{-5}$ Ha eventually & 6--9k calls & 286 calls to gap$<10^{-2}$ (tiny) & consistent \\
Gradient methods (BFGS) struggle with noise & ES/BFGS plateau higher & BFGS/CG never reach gap$<10^{-1}$ & yes \\
ExcitationSolve initially fast then plateaus above surrogate LS & Fig.~2 & not benchmarked & --- \\
40-qubit IBM QPU within 2.5 SD after 3 iters & Fig.~4 & not testable (no IBM access) & --- \\
Depolarizing/coherent gate errors dominate on QPU, $N_s{=}12$--$32$ & Fig.~5 & not testable & --- \\
\bottomrule
\end{longtable}

\section{Verdict: PARTIAL}

\textbf{Rubric application.}
\begin{itemize}
  \item \textbf{REPLICATED} would require reproducing the paper's headline numbers
        ($2$--$4\times$ speedup at $\delta E = 10^{-3}$--$10^{-5}$ Ha on
        H$_2$O/N$_2$/H$_4$, plus the 40-qubit IBM QPU result). This requires the SWS
        simulator (not publicly released) and IBM Quantum access. Not honestly available.
  \item \textbf{PARTIAL} fits: I pulled real public code (STALK v0.1), independently
        reimplemented the paper's Hamiltonian and ansatz from scratch, ran a real
        benchmark against 5 traditional optimizers, and reproduced the qualitative
        direction and rough magnitude ($2.5\times$ vs paper's $2$--$4\times$) of the
        paper's central claim on a scaled-down version of the paper's own TFIM problem.
  \item \textbf{SPOT-CHECK} would apply if I only verified availability. I did more.
  \item \textbf{FAILED} / \textbf{CONTRADICTED} do not fit --- nothing observed
        contradicts the paper; results agree with the paper's direction on every
        testable claim.
  \item \textbf{NO-GO} would apply if the paper's data/code were entirely unavailable;
        STALK v0.1 is available and I used it (in reimplementation).
\end{itemize}

\section{Genuine Critique}

This is a section of substantive critique, distinguishing (a) reproducibility gaps that
belong to the paper, (b) methodological choices that the paper defends but could be
questioned, and (c) limitations of my replication itself.

\subsection{Reproducibility gaps that belong to the paper}

\begin{enumerate}
  \item \textbf{The sparse wave function simulator (SWS) is not released.} The paper's
        chemistry headline numbers (H$_2$O, N$_2$, H$_4$ in Table~1 and Fig.~2) all
        depend on SWS as the surrogate. Refs 86, 87 are cited, but no public URL is
        given in the paper's Data/Materials/Software Availability section, and SWS is
        not part of STALK v0.1. This is a genuine reproducibility gap. A reader cannot
        reproduce the chemistry benchmarks without either (a) writing an independent
        sparse-wavefunction simulator or (b) obtaining SWS from the authors privately.
        For a PNAS paper whose central novelty is a surrogate-based optimizer, the
        surrogate itself being unavailable is a real limitation.
  \item \textbf{The 40-qubit QPU result depends on a paid, gated resource.} The
        \texttt{ibm\_brisbane} demonstration cannot be reproduced without IBM Quantum
        credits. This is not the authors' fault, but it makes Fig.~4 and Fig.~5
        third-party-non-verifiable. The paper does not release the raw QPU shot data
        (bitstring counts per circuit per iteration) that would allow post-hoc
        analysis without re-running on hardware.
  \item \textbf{Baseline optimizer configurations not fully specified.} The paper
        reports that gradient-based methods ``fail'' under sampling noise, but the
        exact finite-difference step size, tolerance settings, and gradient estimation
        method used for BFGS/CG are not stated. This is important because SPSA (which
        is designed for noisy gradients) is not among the baselines --- and its
        exclusion is not justified. A charitable read is that SPSA is well known and
        the paper focuses on scipy defaults; a critical read is that the strongest
        gradient-based baseline for noisy VQE was omitted.
\end{enumerate}

\subsection{Methodological choices worth challenging}

\begin{enumerate}
  \item \textbf{Surrogate quality assumption.} The paper's argument implicitly assumes
        that a small-bond-dimension MPS is a ``good enough'' surrogate for the true
        VQE cost landscape --- specifically, that its Hessian eigenvectors are close
        to the true Hessian's eigenvectors. This is plausible for weakly entangled
        systems but is not obviously true for the strongly entangled regime where VQE
        is most useful. The paper's chemistry tests are all in regimes where MPS/SWS
        works well; the method's failure modes on strongly entangled circuits are not
        characterized.
  \item \textbf{4-parameter TFIM ansatz is intentionally restrictive.} The paper's
        40-qubit demonstration uses a 4-parameter ansatz, which cannot reach the true
        ground state (our replication confirms: at $N_s{=}4$, ansatz min is $-0.686$
        vs.\ true GS $-3.945$). This is fine as an \emph{optimizer benchmark} but
        should not be conflated with a demonstration that VQE + surrogate LS solves
        the physical problem. The paper is careful in language but a casual reader
        could over-interpret the QPU result.
  \item \textbf{No wall-clock or dollar-cost accounting.} Function-call counts are the
        reported metric. But the surrogate itself has non-trivial cost (Hessian
        computation on MPS or SWS), and the paper does not compare total wall-clock
        or total resource cost across methods. On chemistry systems where SWS is
        cheap this may not matter; on larger systems it could dominate.
  \item \textbf{No comparison to Bayesian optimization / RL-based optimizers.}
        Gaussian-process-based Bayesian optimization is a well-established
        alternative for noisy black-box function optimization and has been applied to
        VQE (e.g.\ Iannelli \& Jansen, Nakaji et al.). Its omission from baselines is
        a methodological weakness --- especially since GP-BO is arguably the closest
        cousin to a ``surrogate optimizer.''
\end{enumerate}

\subsection{Limitations of this replication}

\begin{enumerate}
  \item Ns$=4$ vs paper's Ns$=40$; the 4-parameter ansatz cannot reach the true GS,
        so I benchmark distance to the ansatz's variational minimum (fair for
        \emph{optimizer performance} but doesn't test quantum advantage).
  \item Surrogate is exact statevector, not MPS(bond-4). This biases in favor of
        SurrogateLS (``perfect'' surrogate). A more faithful test would use MPS(bond-4)
        as surrogate and exact-statevector-plus-noise as high-level; feasible on
        uicgpu but requires meaningful TN engineering not in scope for this session.
  \item ExcitationSolve --- the paper's tightest competitor to SurrogateLS --- was
        not benchmarked (not in scipy; would require pulling ref 95's implementation).
        The paper's most informative comparison is therefore not independently
        checked.
  \item Chemistry benchmarks (H$_2$O, N$_2$, H$_4$) not attempted, for the SWS
        reasons above.
\end{enumerate}

\subsection{Bottom line}

The paper's methodology is honest and its central qualitative claims survive an
independent partial-scale test. The quantitative headline numbers cannot be
independently verified within a single session without (a) the private SWS simulator
and (b) IBM Quantum access --- a paper-side reproducibility gap, not a replicator-side
failure. A reader should treat the ``$2$--$4\times$ fewer function calls'' claim as
directionally supported at least on TFIM under our simpler-than-paper conditions,
and reserve full judgment on the chemistry headline numbers until either SWS is
released or an independent MPS-surrogate reimplementation is carried out.

\section*{Evidence Files}
\begin{itemize}
  \item \texttt{report/evidence/vqe\_ising\_run.log} --- v1 single-seed run log
  \item \texttt{report/evidence/vqe\_ising\_results.json} --- v1 results JSON
  \item \texttt{report/evidence/vqe\_ising\_v2\_run.log} --- v2 multi-seed run log
  \item \texttt{report/evidence/vqe\_ising\_results\_v2.json} --- v2 results JSON
  \item \texttt{work/replicate\_vqe\_ising.py} --- v1 script
  \item \texttt{work/replicate\_vqe\_ising\_v2.py} --- v2 script
  \item \texttt{work/paper.pdf} --- PNAS PDF (5.0 MB, MD5 df95983\ldots)
  \item \texttt{work/code/stalk-v0.1.tar.gz} --- STALK source (181 KB, MD5 b7e6e41\ldots)
\end{itemize}

\end{document}
